basics of Git and GitHub, including how to create a repository, commit changes, and push to a remote repository.

learning Objectives:
- Understand the basics of Git and GitHub
- Learn how to create a new repository on GitHub
- Learn how to clone a repository to your local machine
- Learn how to make changes to files and commit those changes
- Learn how to push changes to a remote repository on GitHub
- Learn how to pull changes from a remote repository to your local machine
- Learn how to resolve merge conflicts when collaborating with others
- Learn how to use branches to manage different versions of your code
- Learn how to use Git commands in the terminal or command line interface
- Learn how to use GitHub features such as issues, pull requests, and project boards
- Learn how to collaborate with others on GitHub by forking repositories and submitting pull requests
- Learn how to use GitHub Actions for continuous integration and deployment
- Learn how to use GitHub Pages to host static websites directly from a repository
- Learn how to use GitHub's security features, such as branch protection rules and secret management
- Learn how to use GitHub's API to automate tasks and integrate with other tools and services
- Learn how to use GitHub's advanced features, such as code review, project management, and team collaboration tools 
- Learn how to use GitHub's analytics and insights to track repository activity and performance
- Learn how to use GitHub's marketplace to discover and integrate third-party tools and services


Commands:
1- Create a new repository on GitHub:
   - Go to GitHub and log in to your account.
   - Click on the "+" icon in the top right corner and select "New repository".
   - Fill in the repository name, description, and choose whether it will be public or private.
   - Click "Create repository".
2- Clone a repository to your local machine:
   - Open your terminal or command line interface.
   - Navigate to the directory where you want to clone the repository.
   - Use the command `git clone <repository_url>` to clone the repository.
   3- Make changes to files and commit those changes:
   - Open the files you want to edit in your preferred text editor or IDE.
    - Make the necessary changes to the files.
    - Save the changes and return to your terminal.
    - Use the command `git add <file_name>` to stage the changes for commit.
    - Use the command `git commit -m "Your commit message"` to commit the changes with a descriptive message.
4- Push changes to a remote repository on GitHub:
   - Use the command `git push origin <branch_name>` to push your committed changes to the remote repository on GitHub.
5- Pull changes from a remote repository to your local machine:
   - Use the command `git pull origin <branch_name>` to fetch and merge changes from the remote repository to your local machine.
6- Resolve merge conflicts when collaborating with others:
   - If you encounter a merge conflict, open the conflicting files in your text editor.
   - Look for the conflict markers (<<<<<<<, =======, >>>>>>>) and decide how to resolve the conflict.
   - Edit the file to remove the conflict markers and keep the desired changes.
   - Save the file and use `git add <file_name>` to stage the resolved file.
   - Use `git commit` to commit the resolved changes.
7- Use branches to manage different versions of your code:
   - Use the command `git branch <branch_name>` to create a new branch.
   - Use the command `git checkout <branch_name>` to switch to the new branch.
   - Make changes and commit them on the new branch.
   - Use `git push origin <branch_name>` to push the new branch to the remote repository.
8- Use Git commands in the terminal or command line interface:
   - Familiarize yourself with common Git commands such as `git status`, `git log`, `git diff`, and `git reset`.
   - Use these commands to check the status of your repository, view commit history, see changes made to files, and undo changes if necessary.
9- Use GitHub features such as issues, pull requests, and project boards:
   - Create issues to track bugs, feature requests, or tasks related to your project.
    - Use pull requests to propose changes to a repository and collaborate with others on code reviews.
    - Use project boards to organize and manage tasks, track progress, and prioritize work.
10- Collaborate with others on GitHub by forking repositories and submitting pull requests:
    - Fork a repository to create your own copy of the project.
    - Make changes to your forked repository and commit those changes.
    - Submit a pull request to the original repository to propose your changes for review and potential inclusion.

    